\documentclass[../../../include/open-logic-section]{subfiles}

\begin{document}

\olfileid{his}{set}{cantorplane}
\olsection{Cantor ngenani Garis lan Bidang}

Sawetara kahanan ing sakupengé bukti
Schr\"oder-Bernstein gegandhèngan
karo sajarah kang dirembug ing
\olref[his][set][pathology]{sec}.
Élinga yèn ing taun 1877 Cantor
mbuktèkaké yèn cacah titik ing
kothak padha persis karo cacah
titik ing salah siji sisihé.
Ing kéné kita bakal nyuguhaké
upaya buktiné kang kapisan.

Anggepen $\unitline$ garis satuan,
yaiku himpunan titik $[0,1]$.
Anggepen $\unitsquare$ kothak
satuan, yaiku himpunan titik
$\unitline \times \unitline$.
Kanthi notasi iki, Cantor
mbuktèkaké yèn
$\cardeq{\unitline}{\unitsquare}$.
Dhèwèké nulis cathetan
marang Dedekind kang
intiné ngemot argumèn
ing ngisor iki.

\begin{thm}\ollabel{thm:cantorplane}
$\cardeq{\unitline}{\unitsquare}$
\end{thm}

\begin{proof}[Bukti: pérangan kapisan.]
Tetepna $a, b \in \unitline$.
Tulis nganggo notasi binèr:
kanggo saben wilangan kurang
saka siji, pilih ekspansi
kang ora pungkasané kebak
angka siji; kanggo wilangan
siji, pilih ekspansi kang
kabèh angkané siji.
Kanthi mangkono kita olèh
runtutan tanpa wates saka
$0$s lan $1$s, yaiku
$a_1$, $a_2$, \dots,
lan $b_1$, $b_2$,
\dots, saéngga:
\begin{align*}
a &= 0.a_1a_2a_3a_4\dots\\
b &= 0.b_1b_2b_3b_4\dots
\intertext{Saiki delengen fungsi $f \colon \unitsquare \to \unitline$ kang diwènèhaké déning}
f(a, b) & = 0.a_1b_1a_2b_2a_3b_3a_4b_4\dots
\end{align*}
Ekspansi kang diasilaké kanthi
nyelang-nyeling angka uga baku:
angka nolé tanpa wates yèn ora
kaloro input padha karo siji;
yèn kaloroné siji, kabèh
angkané siji. Mula $f$ iku
!!a{injection}, amarga yèn
$f(a, b) = f(c,d)$, mula
$a_n = c_n$ lan $b_n = d_n$
kanggo saben $n \in \Nat$,
saéngga $a = c$ lan $b = d$.
\end{proof}

Sayangé, kaya kang dituduhaké
Dedekind marang Cantor, iki
durung mangsuli pitakonan
wiwitan. Delengen $0.00\dot{1}\dot{0} =
0.0010101010\ldots$.
Supaya $f(a,b) = 0.00\dot{1}\dot{0}$,
angka-angkané bakal meksa:
\begin{align*}
a&= 0.0\dot{1} = 0.011111\ldots\\
b&= 0
\end{align*}
Nanging $a = 0.0\dot{1} = 0.1$.
Mula nalika kandha ``tulis
$a$ lan $b$ nganggo notasi
binèr'', kita kudu netepaké
\emph{ekspansi endi} kang
dienggo; supaya $f$ dadi
\emph{fungsi}, mung
\emph{siji} saka rong
ekspansi kang mungkin
kanggo wilangan iki
sing kena dipilih.
Miturut pilihan baku
ing ndhuwur, $a$
kudu ditulis
``$0.1000\ldots$'',
dudu nganggo angka siji
kang mbaleni. Amarga
output kang dituju
nduwèni ekspansi
binèr kang tunggal,
ora ana pasangan
$\tuple{a, b}$ kang
ngasilaké
$f(a, b) = 0.00\dot{1}\dot{0}$.

Mula, Dedekind nuduhaké
yèn amarga sawetara
ekspansi binèr bisa
mbaleni, fungsi
Cantor~$f$ iku
!!a{injection}
nanging \emph{dudu}
!!a{surjection}.
Dadi Cantor mung
nuduhaké yèn
$\cardle{\unitsquare}{\unitline}$,
\emph{dudu} yèn
$\cardeq{\unitsquare}{\unitline}$.

Cantor meh langsung mbales
layangé Dedekind lan
intiné nyaranaké
supaya bukti mau
dirampungaké mangkéné:

\begin{proof}[Bukti: dirampungaké.]
Dadi kita wis nuduhaké
$\cardle{\unitsquare}{\unitline}$.
Nanging cetha ana
!!a{injection} saka
$\unitline$ menyang
$\unitsquare$: cukup
lebokna garis iku
ing salah siji sisih
kothak. Mula
$\cardle{\unitline}{\unitsquare}$ lan
$\cardle{\unitsquare}{\unitline}$. Miturut Schr\"{o}der--Bernstein
(\olref[sfr][siz][sb]{thm:schroder-bernstein}),
$\cardeq{\unitline}{\unitsquare}$.
\end{proof}

Nanging mesthi waé Cantor ora
bisa ngrampungaké baris
pungkasan nganggo cara iki,
amarga Teorema
Schr\"{o}der-Bernstein
wektu iku durung dibuktèkaké.
Sanajan Cantor mengko
ngrumusaké minangka
konjektur umum,
teorema iki durung
dibuktèkaké kanthi
marem nganti taun 1897.
(Mula, ing pungkasan
taun 1877, Cantor
mènèhi bukti liya
kanggo
\olref{thm:cantorplane}
kang ora lumantar
Schr\"{o}der--Bernstein.)

\end{document}
